\subsection{GERMS WITH RESPECT TO A FILTER}

Let \(\mathfrak{F}\) be a filter on a set \(X\) . On the set \(\mathfrak{P}(X)\) of all subsets of \(X\) , the relation

\[
" \text { there   exists } \quad V \in \mathfrak {F} \text { such   that } \quad M \cap V = N \cap V"
\]

between M and N is an equivalence relation R, for R is obviously reflexive and symmetric, and if M, N, P are three subsets of X such that \(M \cap V = N \cap V\) and \(N \cap W = P \cap W\) , where V and W belong to F, it follows that \(M \cap (V \cap W) = N \cap (V \cap W) = P \cap (V \cap W)\) and \(V \cap W \in F\) , so that R is transitive. The equivalence class mod R of a subset M of X is called the germ of M with respect to F; the quotient set \(\mathfrak{P}(X)/R\) is called the set of germs of subsets of X (with respect to F).

The mappings \((\mathbf{M},\mathbf{N})\to \mathbf{M}\cap \mathbf{N}\) and \((\mathbf{M},\mathbf{N})\rightarrow \mathbf{M}\cup \mathbf{N}\) of

\[
\mathfrak {P} (\mathbf {X}) \times \mathfrak {P} (\mathbf {X})
\]

into \(\mathfrak{P}(\mathbf{X})\) are compatible with the equivalence relations \(\mathbf{R} \times \mathbf{R}\) and \(\mathbf{R}\) (Set Theory, R, § 5, no. 8). For if \(\mathbf{M} \equiv \mathbf{M}'\) (mod R) and \(\mathbf{N} \equiv \mathbf{N}'\) (mod R) then there exist \(\mathbf{V}\) and \(\mathbf{W}\) in \(\mathfrak{F}\) such that

\[
\begin{array}{r l} & {\mathbf {M} \cap \mathbf {V} = \mathbf {M} ^ {\prime} \cap \mathbf {V} \qquad \mathrm{and} \qquad \mathbf {N} \cap \mathbf {W} = \mathbf {N} ^ {\prime} \cap \mathbf {W},} \\ {\mathrm{so\ that}} & {(\mathbf {M} \cap \mathbf {N}) \cap (\mathbf {V} \cap \mathbf {W}) = (\mathbf {M} ^ {\prime} \cap \mathbf {N} ^ {\prime}) \cap (\mathbf {V} \cap \mathbf {W})} \\ {\mathrm{and}} & {(\mathbf {M} \cup \mathbf {N}) \cap (\mathbf {V} \cap \mathbf {W}) = (\mathbf {M} \cap (\mathbf {V} \cap \mathbf {W})) \cup (\mathbf {N} \cap (\mathbf {V} \cap \mathbf {W}))} \\ & {\qquad = (\mathbf {M} ^ {\prime} \cap (\mathbf {V} \cap \mathbf {W})) \cup (\mathbf {N} ^ {\prime} \cap (\mathbf {V} \cap \mathbf {W}))} \\ & {\qquad = (\mathbf {M} ^ {\prime} \cup \mathbf {N} ^ {\prime}) \cap (\mathbf {V} \cap \mathbf {W}).} \end{array}
\]

Passing to the quotients, these mappings induce two mappings of \((\mathfrak{P}(\mathbf{X}) / \mathbf{R}) \times (\mathfrak{P}(\mathbf{X}) / \mathbf{R})\) into \(\mathfrak{P}(\mathbf{X}) / \mathbf{R}\) , which we denote (by abuse of language) by \((\xi, \eta) \to \xi \cap \eta\) and \((\xi, \eta) \to \xi \cup \eta\) respectively. It is a straightforward exercise to verify that with respect to these laws of composition {[}defined throughout \(\mathfrak{P}(\mathbf{X}) / \mathbf{R}\) {]} every element is idempotent, and that each law is commutative and associative and distributive with respect to the other. Further, the relations \(\xi = \xi \cap \eta\) and \(\eta = \xi \cup \eta\) are equivalent; if we denote them (by abuse of language) by \(\xi \subset \eta\) , it is easily verified that this relation is an ordering on \(\mathfrak{P}(\mathbf{X}) / \mathbf{R}\) , with respect to which \(\mathfrak{P}(\mathbf{X}) / \mathbf{R}\) is a lattice which has the germ of \(\varnothing\) as least element and the germ of \(\mathbf{X}\) as greatest element. Note that the relation \(\xi \subset \eta\) means that there exist \(\mathbf{M} \in \xi\) , \(\mathbf{N} \in \eta\) and \(\mathbf{V} \in \mathfrak{F}\) such that \(\mathbf{M} \cap \mathbf{V} \subset \mathbf{N} \cap \mathbf{V}\) .

Now let \(X'\) be another set, and let \(\Phi\) be the set of all mappings of a set of F into \(X'\) . The relation on \(\Phi\)

``there exists \(V \in F\) such that f and g are defined and agree on \(V\)''

between \(f\) and \(g\) is an equivalence relation \(\mathbf{S}\) ; it is clear that \(\mathbf{S}\) is reflexive and symmetric, and if \(f, g, h\) are three elements of \(\Phi\) such that \(f\) and \(g\) are defined on \(\mathbf{V} \in \mathfrak{F}\) , and \(g\) and \(h\) are defined and agree on \(\mathbf{W} \in \mathfrak{F}\) , then \(f\) and \(h\) are defined and agree on \(\mathbf{V} \cap \mathbf{W} \in \mathfrak{F}\) , so that \(\mathbf{S}\) is transitive. The equivalence class mod \(\mathbf{S}\) of a mapping \(f\) of a set \(\mathbf{V} \in \mathfrak{F}\) into \(\mathbf{X}'\) is called the germ of \(f\) (with respect to \(\mathfrak{F}\) ), and the quotient set \(\tilde{\Phi} = \Phi / \mathbf{S}\) is called the set of germs of mappings of \(\mathbf{X}\) into \(\mathbf{X}'\) (with respect to \(\mathfrak{F}\) ).

Remarks. 1) Every mapping \(f\) of a subset \(M\) of \(X\) into \(X'\) , where \(M\) belongs to \(\mathfrak{F}\) , is equivalent mod \(S\) to a mapping \(f_1\) of \(X\) into \(X'\) (which justifies the above terminology): it is sufficient to extend \(f\) to \(X\) , e.g.~by giving it a constant value on \(X - M\) .

\begin{enumerate}
\def\labelenumi{\arabic{enumi})}
\setcounter{enumi}{1}
\tightlist
\item
  The characteristic functions \(\varphi_{\mathbf{M}}\) and \(\varphi_{\mathbf{N}}\) of two subsets \(\mathbf{M}\) and \(\mathbf{N}\) of \(\mathbf{X}\) have the same germ with respect to \(\mathfrak{F}\) if and only if \(\mathbf{M}\) and \(\mathbf{N}\) have the same germ with respect to \(\mathfrak{F}\) .
\end{enumerate}

Let \(X''\) be a third set, \(\varphi\) a mapping of \(X'\) into \(X''\) , \(\Phi'\) the set of all mappings of a set of \(\mathfrak{F}\) into \(X''\) . For each \(f \in \Phi\) , \(\varphi \circ f\) belongs to \(\Phi'\) ; further, it is immediately seen that if \(g \in \Phi\) has the same germ as \(f\) with respect to \(\mathfrak{F}\) , then \(\varphi \circ f\) and \(\varphi \circ g\) have the same germ with respect to \(\mathfrak{F}\) .

This germ therefore depends only on the germ \(\tilde{f}\) of \(f\) with respect to \(\mathfrak{F}\) and is denoted by \(\varphi(\tilde{f})\) . We thus define a mapping (denoted by \(\varphi\) , by abuse of language) of the set \(\tilde{\Phi}\) of germs of mappings of X into \(\mathrm{X}'\) , into the set \(\tilde{\Phi}'\) of germs of mappings of X into \(\mathrm{X}''\) .

Now let \(\mathbf{X}_i^{\prime}\) ( \(1 \leqslant i \leqslant n\) ) be sets and

\[
\mathbf {Y} = \prod_ {i = 1} ^ {n} \mathbf {X} _ {i} ^ {\prime}
\]

their product; let \(\Phi_i\) (resp. \(\Phi\) ) denote the set of all mappings of a set of \(\mathfrak{F}\) into \(\mathbf{X}_i'\) (resp. \(\mathbf{Y}\) ). If \(f_i \in \Phi_i\) for \(1 \leqslant i \leqslant n\) and if \(\mathbf{M}_i \in \mathfrak{F}\) is the domain of \(f_i\) , then the mapping \(t \to (f_1(t), \ldots, f_n(t))\) is defined on

\[
\bigcap_ {i = 1} ^ {n} \mathbf {M} _ {i}
\]

and hence belongs to \(\Phi\) ; we denote this mapping (by abuse of language) by \((f_1, \ldots, f_n)\) . Furthermore, if \(f_i\) and \(g_i\) belong to \(\Phi_i\) and have the same germ with respect to \(\mathfrak{F}\) (for \(1 \leqslant i \leqslant n\) ), it is immediately seen that \((f_1, \ldots, f_n)\) and \((g_1, \ldots, g_n)\) have the same germ with respect to \(\mathfrak{F}\) ; this germ therefore depends only on the germs \(\tilde{f}_i\) of the \(f_i\) . If we denote it by \(\Gamma(\tilde{f}_1, \ldots, \tilde{f}_n)\) then \(\Gamma\) is clearly a bijection of the product set

\[
\prod_ {i = 1} ^ {n} \tilde {\Phi} _ {i}
\]

onto the set \(\tilde{\Phi}\) , where \(\tilde{\Phi}_i\) (resp. \(\tilde{\Phi}\) ) denotes the set of germs of mappings of X into \(\mathbf{X}_i'\) (resp. Y) with respect to \(\mathfrak{F}\) ; hence, by abuse of language, we shall generally write \((\tilde{f}_1, \ldots, \tilde{f}_n)\) instead of \(\Gamma(\tilde{f}_1, \ldots, \tilde{f}_n)\) whenever there is no risk of confusion.

From what has been said, every mapping \(\psi\) of Y into a set \(\mathbf{X}''\) defines a mapping \((\tilde{f}_1, \ldots, \tilde{f}_n) \to \psi(\tilde{f}_1, \ldots, \tilde{f}_n)\) of

\[
\prod_ {i = 1} ^ {n} \tilde {\Phi} _ {i}
\]

into the set \(\tilde{\Phi}'\) of germs of mappings of X into \(\mathbf{X}''\) .

In particular, if \(n = 2\) and if \(X_1', X_2'\) and \(X''\) are all equal to the same set \(X'\) (so that \(\psi\) is a law of composition defined throughout \(X'\) ), then \(\psi\) induces a law of composition defined throughout the set \(\tilde{\Phi}\) of germs of mappings of \(X\) into \(X'\) . It is easily verified that if the law given on \(X'\) is associative (resp. commutative) then so is the corresponding law on \(\tilde{\Phi}\) ; if the law \(\psi\) on \(X'\) has an identity element \(e'\) , then the germ with respect to \(\mathfrak{F}\) of the constant mapping \(x \to e'\) is an identity element for the corresponding law on \(\tilde{\Phi}\) . Finally, if the law on \(X'\) has an identity element \(e'\) , then the germ \(\tilde{f}\) of \(f \in \Phi\) has an inverse in \(\tilde{\Phi}\) if and only if there exists \(V \in F\) , contained in the domain of f, such that \(f(t)\) is invertible in \(X'\) for each \(t \in V\) ; if, for each \(t \in V\) , \(g(t)\) denotes the inverse of \(f(t)\) then the germ \(\tilde{g}\) of g is the inverse of \(\tilde{f}\) in \(\tilde{\Phi}\) . In particular if \(X'\) is a group with respect to the law \(\psi\) , then \(\tilde{\Phi}\) is a group with respect to the corresponding law; likewise, if \(X'\) is a ring (resp. an algebra over a ring A) then \(\tilde{\Phi}\) is a ring (resp. an algebra over A) with respect to the corresponding laws of composition.

